\chapter{Относительный кинетический вес из общего коэффициента теплового ядра}
\label{ch:v8-full-field-kinetic-relative-weight-parent-origin-gate}

Предыдущая сборка оставила независимые веса скалярного и калибровочного
блоков. Старый четырёхмерный дираковский подъём предлагает проверить единый
оператор
\begin{equation}
 \mathcal D=i\gamma^\mu(\partial_\mu+B_\mu)+\gamma^5\Phi,
 \qquad \{\gamma^\mu,\gamma^\nu\}=2\delta^{\mu\nu}I_4.
 \label{eq:v8-common-dirac-kinetic-lift}
\end{equation}
В плоском фоне локальный коэффициент $a_4$ содержит
\begin{equation}
 (4\pi)^{-2}\Tr\left(\frac12E^2+\frac1{12}\Omega_{\mu\nu}\Omega^{\mu\nu}\right).
 \label{eq:v8-common-a4-kinetic-parent}
\end{equation}
Точный след по явному четырёхмерному реперу Клиффорда даёт
\begin{equation}
 c_{\Phi}=2,
 \qquad c_F=-\frac23
 \quad\text{для антиэрмитовой кривизны},
 \qquad c_F^{+}=\frac23.
 \label{eq:v8-a4-scalar-gauge-coefficients}
\end{equation}
Поэтому общий след на полном внутреннем репере фиксирует
\begin{equation}
 \frac{w_{\text{п}}}{w_{\text{к}}}
 =\frac{c_{\Phi}}{c_F^{+}}=3,
 \qquad (w_{\text{п}},w_{\text{к}})\sim\left(1,\frac13\right).
 \label{eq:v8-a4-relative-kinetic-weight}
\end{equation}

\begin{theorem}[Условный выбор относительного веса]
В стандартном четырёхмерном лапласовом подъёме одного оператора Дирака
общий коэффициент $a_4$ устраняет независимость двух кинетических весов и
точно выбирает отношение $3:1$. Общий момент спектральной функции умножает
оба блока одинаково и отношения не меняет.
\end{theorem}

\begin{nogocriterion}[Конечный родитель ещё не выводит дираковский подъём]
Полученное отношение нельзя объявлять безусловным следствием одного
$42$-мерного конечного репера. Оно использует дополнительную структуру:
четырёхмерный клиффордов модуль, лапласов квадрат и выбор именно общего
$a_4$-действия. Пока этот пространственно-временной подъём не выведен из
текущего родителя, результат имеет статус строгого условного кандидата.
\end{nogocriterion}

\begin{protocol}
\begin{enumerate}
 \item клиффордовы тождества: \textbf{$16$ точных проверок};
 \item скалярный коэффициент: \textbf{$2$};
 \item положительный калибровочный коэффициент: \textbf{$2/3$};
 \item отношение весов: \textbf{$3:1$};
 \item единый внутренний след: \textbf{совместим};
 \item происхождение дираковского подъёма: \textbf{не выведено};
 \item безусловный физический вес: \textbf{не закрыт};
 \item следующий гейт: \textbf{происхождение общего $a_4$-дираковского подъёма}.
\end{enumerate}
\end{protocol}

Машинный сертификат сохранён в
\path{s2t_v8_full_field_kinetic_relative_weight_parent_origin_gate_results.json}.